\documentclass[11pt]{article}
\usepackage[margin=1.2in]{geometry}
\usepackage{amsmath,amssymb,amsthm}
\usepackage{enumitem}
\usepackage{hyperref}
\usepackage{booktabs}

\newcommand{\D}{\mathbb{D}}
\newcommand{\R}{\mathbb{R}}
\newcommand{\C}{\mathbb{C}}
\newcommand{\Z}{\mathbb{Z}}
\newcommand{\rhoT}{\rho_T}
\newcommand{\dzbar}{\partial_{\bar{z}_i}}

\theoremstyle{plain}
\newtheorem{conjecture}{Conjecture}
\newtheorem{question}{Question}
\newtheorem{observation}{Observation}
\theoremstyle{definition}
\newtheorem{definition}{Definition}
\theoremstyle{remark}
\newtheorem{remark}{Remark}

\title{Combinatorics of the Poincar\'e Transport Residual:\\
  Discretisation, Monodromy, and Cylindric Partitions}
\author{Speculative notes from a conversation}
\date{March 2026}

\begin{document}
\maketitle

\begin{abstract}
We sketch a network of connections between the Poincar\'e transport
residual~$\rhoT$ of a viscous elliptic structure (VES) and the
combinatorics of exactly solvable lattice models. The central
observation is that the rigidity condition~$G=0$ is a continuous
analogue of the Yang--Baxter equation, and that~$\rhoT$ measures
distance from exact solvability. When discretised on a lattice, the
VES compatibility condition yields local rules of the type appearing
in growth diagrams, puzzle pieces, and transfer-matrix formulations of
Schur functions. We propose that the profile of a cylindric partition
is a discrete monodromy, the combinatorial counterpart of~$\rhoT$.
\end{abstract}

\tableofcontents

%% ────────────────────────────────────────────────────
\section{The Poincar\'e transport residual}
%% ────────────────────────────────────────────────────

Let~$\Omega\subset\R^2$ be a domain. A \emph{viscous elliptic
structure} (VES) on~$\Omega$ is specified by a smooth generator
$i\colon\Omega\to\D$, where~$\D$ is the open unit disk. The
associated Cauchy--Riemann operator is
\[
  \dzbar f \;=\; \tfrac{1}{2}\bigl(f_x + i(x,y)\,f_y\bigr),
\]
and a function~$f$ is \emph{$i$-holomorphic} when $\dzbar f = 0$.

\begin{definition}[Rigidity obstruction]
The \emph{obstruction tensor} of a VES is
\[
  G \;=\; i_x + i\cdot i_y.
\]
The structure is \emph{rigid} (admits a full holomorphic calculus)
if and only if $G=0$.
\end{definition}

\begin{definition}[Poincar\'e transport residual]
The \emph{Poincar\'e transport residual} is
\[
  \boxed{\rhoT \;=\; \frac{|G|}{1-|i|^2}
         \;=\; \frac{|i_x + i\cdot i_y|}{1-|i|^2},}
\]
where $(1-|i|^2)^{-1}$ is the conformal factor of the Poincar\'e
metric at $i\in\D$.
\end{definition}

Key properties:
\begin{enumerate}[label=(\roman*)]
  \item \textbf{Invariance.} $\rhoT$ is invariant under M\"obius
    automorphisms of~$\D$ acting on~$i$---it is a genuine geometric
    invariant.
  \item \textbf{Blow-up.} As $|i|\to 1$ (degenerate ellipticity),
    $\rhoT\to\infty$ even for bounded~$G$, correctly detecting
    the parabolic/hyperbolic boundary.
  \item \textbf{Vanishing.} $\rhoT=0$ if and only if $G=0$
    (rigidity).
  \item \textbf{First-order.} $\rhoT$ involves one derivative
    of~$i$; it detects how the structure \emph{varies}, not just
    what it is at a point.
\end{enumerate}

\begin{remark}[The difficulty plane]
The pair $(\rhoT,\,|\mu|)$---where $|\mu|$ is the Beltrami
modulus---defines a two-dimensional ``difficulty plane.'' The
invariant~$|\mu|$ measures algebraic deviation (anisotropy);
$\rhoT$ measures differential deviation (spatial variation). They
are categorically independent.
\end{remark}


%% ────────────────────────────────────────────────────
\section{The Burgers connection}
%% ────────────────────────────────────────────────────

The obstruction $G = i_x + i\cdot i_y$ has the algebraic form of the
\emph{Burgers nonlinearity}~$u\,u_x$ in the inviscid Burgers equation
$u_t + u\,u_x = 0$. This is not a coincidence: the VES monograph
identifies~$G$ as the nonlinear term in the Burgers transport equation
for the spectral parameter~$\lambda$.

The Burgers equation is the gateway to a remarkable chain of
connections.

\subsection{The chain: Burgers $\to$ KPZ $\to$ Schur}

\begin{enumerate}[label=\textbf{Step \arabic*.}, leftmargin=2.5em]

\item \textbf{Burgers $\to$ KPZ.}
The Hopf--Cole transformation $u = -2\nu\,\partial_x\log Z$
linearises Burgers into the heat equation. Adding space-time white
noise~$\eta$ gives the \emph{KPZ equation}
(Kardar--Parisi--Zhang, 1986):
\[
  \partial_t h = \nu\,\partial_x^2 h
    + \tfrac{\lambda}{2}(\partial_x h)^2 + \eta(x,t).
\]
Setting $u=\lambda\,\partial_x h$ recovers stochastic Burgers; the
Hopf--Cole transform $Z=\exp(\lambda h/2\nu)$ gives the stochastic
heat equation with multiplicative noise, whose solution is a
\emph{directed polymer} partition function.

\item \textbf{KPZ $\to$ growth models.}
KPZ is the continuum limit of discrete growth models. The cleanest
is \emph{last passage percolation} (LPP): given i.i.d.\ weights
$w(i,j)$ on~$\Z_{>0}^2$, define
\[
  G(m,n) = \max_{\pi:(1,1)\to(m,n)}
    \sum_{(i,j)\in\pi} w(i,j)
\]
over up-right lattice paths~$\pi$. The surface~$h(m,n)=G(m,n)$
grows in the KPZ universality class with fluctuation
exponent~$\beta=1/3$.

The \emph{totally asymmetric simple exclusion process} (TASEP)---
particles on~$\Z$ hopping right at rate~1 if the target is
empty---is equivalent to LPP (Rost, 1981).

\item \textbf{Growth models $\to$ Schur functions (via RSK).}
The Robinson--Schensted--Knuth correspondence bijects
nonneg-integer matrices to pairs of semistandard Young tableaux.
Greene's theorem (1974): the shape~$\lambda_1$ equals the longest
weakly increasing subsequence.

For LPP with geometric weights $w(i,j)\sim\mathrm{Geom}(x_i y_j)$:
the last passage time~$G(m,n)$ equals~$\lambda_1$ under RSK.
The random partition is distributed according to the \emph{Schur
measure} (Okounkov, 2001):
\[
  \mathrm{Prob}(\lambda) \;\propto\;
    s_\lambda(x_1,\ldots,x_m)\,s_\lambda(y_1,\ldots,y_n),
\]
derived from the Cauchy identity
$\sum_\lambda s_\lambda(x)\,s_\lambda(y) =
  \prod_{i,j}(1-x_iy_j)^{-1}$.

\item \textbf{Schur measure $\to$ Tracy--Widom.}
The Schur measure is a \emph{determinantal point process}. The
probability $\mathrm{Prob}(\lambda_1\le t)$ is a Fredholm
determinant; steepest-descent analysis yields the \emph{Airy
kernel}, and
\[
  F_2(s) = \det(I-K_{\mathrm{Airy}})_{L^2(s,\infty)}
         = \exp\!\Bigl(-\int_s^\infty (x-s)\,q(x)^2\,dx\Bigr),
\]
where $q$~solves Painlev\'e~II: $q''=sq+2q^3$. This is the
\emph{Tracy--Widom distribution}, first discovered as the
fluctuation law for the largest eigenvalue of GUE random
matrices (Tracy--Widom, 1994).

\end{enumerate}

\subsection{The role of Yang--Baxter and transfer matrices}

The exact solvability at every level of this chain rests on the
Yang--Baxter equation (YBE). For a vertex model with weights
$R_{ab}^{cd}(u)$, the YBE ensures that transfer matrices
$T(u) = \mathrm{tr}_{\mathrm{aux}}\prod_i R_i(u)$ commute:
$[T(u),T(v)]=0$. This is why:
\begin{itemize}
  \item Schur functions are determinants (Jacobi--Trudi =
    Lindstr\"om--Gessel--Viennot for non-intersecting paths).
  \item The Schur measure is determinantal (commuting transfer
    matrices $\Rightarrow$ free fermions).
  \item The Fredholm determinant is computable in closed form.
\end{itemize}

The transfer operators $T_{\mathrm{free}}$ and
$\bar{T}_{\mathrm{free}}$ of Zinn-Justin~\cite{ZJ2008}, acting on
the fermionic Fock space of binary strings, sit precisely in this
framework. The Littlewood--Richardson coefficients computed via
puzzle pieces control the product structure of Schur functions,
governing multi-point correlations in the Schur process.

%% ────────────────────────────────────────────────────
\section{Discretisation: from $\rhoT$ to local rules}
\label{sec:discretisation}
%% ────────────────────────────────────────────────────

\begin{observation}
Discretise~$\Omega$ on a grid. At each vertex $(x,y)$, place a
generator $i(x,y)\in\D$. The rigidity condition $G=i_x+i\cdot i_y=0$
becomes a \emph{nearest-neighbour constraint}: the values of~$i$ at
adjacent sites must satisfy a compatibility relation.
\end{observation}

This is exactly the setup for:
\begin{itemize}
  \item \textbf{Puzzle pieces} (Knutson--Tao--Woodward): local rules
    constraining adjacent tiles.
  \item \textbf{Growth diagrams} (Fomin): local rules propagating
    partitions through a grid, where the local rule is a
    zero-curvature condition around each plaquette.
  \item \textbf{Vertex models} (six-vertex, ice model): local weights
    depending on adjacent edge states.
\end{itemize}

The generator~$i\in\D$ plays the role of the local state; $G=0$
plays the role of the local rule. The fermionic Fock space of binary
strings (encoding partitions) is a discretisation of the space of
generators.

\begin{question}
Is there a natural discretisation of $G=0$ on a triangular or
square lattice that recovers the interlacing condition on
partitions? If so, $\rhoT$ would literally measure \emph{how much
interlacing fails}.
\end{question}

\subsection{The difficulty plane as (inhomogeneity, anisotropy)}

In lattice model language, the two axes of the difficulty plane map
to fundamental distinctions:

\medskip
\begin{center}
\begin{tabular}{@{}lll@{}}
\toprule
\textbf{VES invariant} & \textbf{Lattice model} & \textbf{Combinatorial effect} \\
\midrule
$|\mu|$ (Beltrami) & Anisotropy of local weights &
  Deviation from isotropy \\
$\rhoT$ (transport) & Inhomogeneity (weights vary site-to-site) &
  Richer combinatorics \\
\bottomrule
\end{tabular}
\end{center}
\medskip

Homogeneous isotropic models give product-formula partition functions.
Homogeneous anisotropic models give the six-vertex model, Schur
functions, determinantal formulas. Inhomogeneous models are harder but
produce the richest combinatorics: factorial Schur functions, double
Schubert polynomials.


%% ────────────────────────────────────────────────────
\section{Profile as discrete monodromy}
\label{sec:monodromy}
%% ────────────────────────────────────────────────────

This is the most speculative---and perhaps most suggestive---
connection.

\subsection{The analogy}

In VES, $\rhoT$ measures the transport obstruction: the monodromy of
the complex structure under parallel transport. The condition~$G=0$
means trivial monodromy.

A cylindric partition on a cylinder of circumference~$d$ with profile
$c=(c_0,c_1,\ldots,c_{r-1})$ has a \emph{discrete monodromy}: as one
traverses the cylinder, the boundary conditions shift by the profile.

\medskip
\begin{center}
\begin{tabular}{@{}ll@{}}
\toprule
\textbf{VES (continuous)} & \textbf{Cylindric partitions (discrete)} \\
\midrule
$G=0$, trivial monodromy & Trivial profile $\to$ plane partitions \\
$G\ne 0$, $\rhoT>0$ & Nontrivial profile $\to$ cylindric partitions \\
$\rhoT$ measures twist magnitude & Profile measures boundary shift \\
\bottomrule
\end{tabular}
\end{center}
\medskip

\begin{conjecture}[Informal]
There exists a discretisation scheme under which the cylindric
profile~$c$ arises as the holonomy of the discrete VES connection
around the cylinder, and~$\rhoT$ (integrated around the cycle)
recovers a function of~$c$.
\end{conjecture}

\subsection{Positivity as ``mild deformation''}

Under this analogy, Warnaar's positivity conjecture---that
$Q_{n,c}(q)\ge 0$ coefficient-wise for rank-3 cylindric
partitions---asserts that \emph{even when $\rhoT>0$, the partition
function retains positivity}. The deformation from exact solvability
is mild enough that combinatorial structure survives without
catastrophic cancellation.

\subsection{The $d\equiv 0\pmod{3}$ resonance}

For rank~3, when $d\equiv 0\pmod{3}$ the profile distributes evenly
among the three parts: the slack vector is
$(d/3,\,d/3,\,d/3)$, which has no cyclic descent. The
cyclic-descent injection rule (which works for all other~$d$) fails
here.

In the continuous picture, this corresponds to \emph{resonant
monodromy}: the monodromy divides evenly into the rank, creating a
fixed point under the cyclic symmetry. In KAM theory, resonant
frequencies are exactly where small-denominator problems arise and
invariant tori break down. The $d\equiv 0\pmod{3}$ obstruction may
be a discrete manifestation of the same phenomenon.


%% ────────────────────────────────────────────────────
\section{Poincar\'e factors and $q$-series}
%% ────────────────────────────────────────────────────

A more speculative thread. The Poincar\'e conformal factor at
$z\in\D$ is $(1-|z|^2)^{-1}$. Imagine a sequence of generators at
different scales with $|i_k|=q^k$. The product of Poincar\'e
factors gives
\[
  \prod_{k=1}^{n}\frac{1}{1-q^{2k}},
\]
which is a partition generating function (partitions into even parts
$\le 2n$). The full $q$-Pochhammer symbol
$(q;q)_\infty = \prod_{n\ge 1}(1-q^n)$---ubiquitous in partition
theory---arises from a product over all levels.

The Euler--Heine factorisation $\mathrm{GK}_c = H(z)\cdot F_c$,
where $H=1/(zq;q)_\infty$, decomposes the generating function into a
``universal'' part (product of Poincar\'e-like factors) and a
``structural'' part~$F_c$. In the VES picture, this might correspond
to:
\[
  \text{universal hyperbolic metric contribution}
    \;\times\;
  \text{structure-specific transport geometry.}
\]

\begin{question}
Is there a multi-level VES (a tower of generators $i_1,i_2,\ldots$
with prescribed nesting) whose ``total Poincar\'e factor'' is the
$q$-Pochhammer symbol? If so, the E--H factorisation would have a
geometric interpretation as separation of metric and transport
contributions.
\end{question}


%% ────────────────────────────────────────────────────
\section{Open directions}
%% ────────────────────────────────────────────────────

\begin{enumerate}
  \item \textbf{Construct the discretisation.} Make
    Section~\ref{sec:discretisation} precise: discretise the VES
    rigidity condition on a square lattice and identify the resulting
    local rule. Does it give interlacing? A vertex model?

  \item \textbf{Deformed Yang--Baxter.} If $\rhoT$ measures distance
    from exact solvability, what do Schur measures look like when
    Yang--Baxter is only approximately satisfied? Is there a
    ``$\rhoT$-deformed'' determinantal structure?

  \item \textbf{Profile from holonomy.} Make the
    Conjecture~in Section~\ref{sec:monodromy} precise: construct a
    discrete VES on a cylindrical lattice whose holonomy recovers the
    cylindric profile~$c$.

  \item \textbf{Resonance and positivity.} Can KAM-type arguments
    (small denominators at resonant frequencies) explain why
    $d\equiv 0\pmod{3}$ is harder for the positivity conjecture?

  \item \textbf{Growth diagrams as flat connections.} Fomin's local
    rule for growth diagrams is a zero-curvature condition around each
    plaquette. Formalise the correspondence: local rule = discrete
    $G=0$, and the failure $G\ne 0$ = ``curvature'' of the growth
    diagram.
\end{enumerate}


%% ────────────────────────────────────────────────────
%% References
%% ────────────────────────────────────────────────────

\begin{thebibliography}{99}

\bibitem{BDJ1999}
J.~Baik, P.~Deift, K.~Johansson,
\emph{On the distribution of the length of the longest increasing
  subsequence of random permutations},
J.~Amer.\ Math.\ Soc.\ \textbf{12} (1999), 1119--1178.

\bibitem{Corwin2012}
I.~Corwin,
\emph{The Kardar--Parisi--Zhang equation and universality class},
Random Matrices Theory Appl.\ \textbf{1} (2012), 1130001.

\bibitem{Fomin1986}
S.~Fomin,
\emph{Generalized Robinson--Schensted--Knuth correspondence},
Zap.\ Nauchn.\ Sem.\ LOMI \textbf{155} (1986), 156--175.

\bibitem{Greene1974}
C.~Greene,
\emph{An extension of Schensted's theorem},
Adv.\ Math.\ \textbf{14} (1974), 254--265.

\bibitem{Johansson2000}
K.~Johansson,
\emph{Shape fluctuations and random matrices},
Comm.\ Math.\ Phys.\ \textbf{209} (2000), 437--476.

\bibitem{KPZ1986}
M.~Kardar, G.~Parisi, Y.-C.~Zhang,
\emph{Dynamic scaling of growing interfaces},
Phys.\ Rev.\ Lett.\ \textbf{56} (1986), 889--892.

\bibitem{Okounkov2001}
A.~Okounkov,
\emph{Infinite wedge and random partitions},
Selecta Math.\ \textbf{7} (2001), 57--81.

\bibitem{OR2003}
A.~Okounkov, N.~Reshetikhin,
\emph{Correlation function of Schur process with application to local
  geometry of a random 3-dimensional Young diagram},
J.~Amer.\ Math.\ Soc.\ \textbf{16} (2003), 581--603.

\bibitem{Romik2015}
D.~Romik,
\emph{The Surprising Mathematics of Longest Increasing Subsequences},
Cambridge Univ.\ Press, 2015.

\bibitem{TW1994}
C.~Tracy, H.~Widom,
\emph{Level-spacing distributions and the Airy kernel},
Comm.\ Math.\ Phys.\ \textbf{159} (1994), 151--174.

\bibitem{ZJ2008}
P.~Zinn-Justin,
\emph{Littlewood--Richardson coefficients and integrable tilings},
Electron.\ J.\ Combin.\ \textbf{16} (2009), R12.
arXiv:0809.2392.

\end{thebibliography}

\end{document}
